\section*{[2023 EXAM] Problem 1: Poisson, Weak form, well-posed}
We are given $u\in H^2(\Omega)$ solving the Poisson problem
\[
  \begin{cases}
    - \Delta u = f & \text{in } \Omega \subset \mathbb{R}^3, \\
    u = 0          & \text{on } \partial\Omega,
  \end{cases}
  \tag{\text{strong}}
\]
for some $f\in L^2(\Omega)$.
Throughout, $\langle\cdot,\cdot\rangle$ denotes the $L^2(\Omega)$ inner product
\[
  \langle u,v\rangle := \displaystyle\int_\Omega u\,v\,dx,
\]
and $\langle\cdot,\cdot\rangle_V$ denotes the inner product of a general Hilbert space $V$.
\subsection*{(a) Definition of $L^2(\Omega)$ and $H^2(\Omega)$}
We first recall the notion of weak derivative.
\textbf{Weak derivative.}
Let $\alpha=(\alpha_1,\alpha_2,\alpha_3)\in\mathbb{N}_0^3$ be a multi--index with
$|\alpha| := \alpha_1+\alpha_2+\alpha_3$.
For $u\in L^1_{\mathrm{loc}}(\Omega)$, a function $w\in L^1_{\mathrm{loc}}(\Omega)$ is called
the \emph{weak derivative} $D^\alpha u$ of order $\alpha$ if
\[
  \int_\Omega u\, D^\alpha\varphi\,dx
  = (-1)^{|\alpha|}\int_\Omega w\,\varphi\,dx
  \quad\forall\varphi\in C_c^\infty(\Omega).
\]
If such a $w$ exists, it is unique (up to sets of measure zero).

\textbf{$L^2(\Omega)$:}
The space $L^2(\Omega)$ is
\[
  L^2(\Omega)
  := \left\{ u : \Omega\to\mathbb{R} \text{ measurable } :
  \int_\Omega |u(x)|^2\,dx < \infty \right\},
\]
equipped with the norm
\[
  \|u\|_{L^2(\Omega)} := \left(\int_\Omega |u|^2\,dx\right)^{1/2}
\]
and inner product $\langle u,v\rangle = \int_\Omega u v\,dx$.

\textbf{$H^2(\Omega)$:}
The Sobolev space $H^2(\Omega)$ is
\[
  H^2(\Omega)
  := \Bigl\{
  u \in L^2(\Omega) \;:\;
  D^\alpha u \in L^2(\Omega)
  \text{ for all multi--indices } \alpha \text{ with } |\alpha|\le 2
  \Bigr\},
\]
with norm
\[
  \|u\|_{H^2(\Omega)}^2
  := \sum_{|\alpha|\le 2} \|D^\alpha u\|_{L^2(\Omega)}^2.
\]
Thus $H^2(\Omega)$ consists of all $L^2$--functions whose weak derivatives up to order $2$
also lie in $L^2(\Omega)$.
\subsection*{(b) Weak formulation: spaces, bilinear form and linear functional}
We now derive the weak formulation of (strong) of the abstract form
\[
  a(u,v) = g(v)
  \quad\forall v\in V,
  \tag{2}
\]
by identifying a Hilbert space $V$, a bilinear form $a(\cdot,\cdot)$ and a bounded linear
functional $g$.

\textbf{Step 1: Multiply by a test function and integrate.}
Let $v\in C_c^\infty(\Omega)$ (later we extend to $H^1_0(\Omega)$ by density).
Multiply (strong) by $v$ and integrate over $\Omega$:
\[
  \int_\Omega (-\Delta u)\,v\,dx
  = \int_\Omega f\,v\,dx.
\]
\textbf{Step 2: Integration by parts and boundary condition.}
Using integration by parts (Green's identity),
\[
  \int_\Omega (-\Delta u)\,v\,dx
  = \int_\Omega \nabla u\cdot\nabla v\,dx
  - \int_{\partial\Omega} \frac{\partial u}{\partial n}\,v\,ds.
\]
Here $\dfrac{\partial u}{\partial n} = \nabla u\cdot n$ is the normal derivative.
Because $u=0$ on $\partial\Omega$, we may choose test functions $v$ that also
vanish on $\partial\Omega$, that is $v\in H^1_0(\Omega)$; then the boundary integral is zero:
\[
  \int_\Omega (-\Delta u)\,v\,dx
  = \int_\Omega \nabla u\cdot\nabla v\,dx
  \quad\forall v\in H^1_0(\Omega).
\]
Thus the weak form is:
\[
  \int_\Omega \nabla u\cdot\nabla v\,dx
  = \int_\Omega f\,v\,dx
  \quad\forall v\in H^1_0(\Omega).
\]
\textbf{Step 3: Identify $V$, $a(\cdot,\cdot)$ and $g$.}
We choose
\[
  V := H^1_0(\Omega),
\]
with inner product
\[
  \langle v,w\rangle_V := \int_\Omega \nabla v\cdot\nabla w\,dx,
  \qquad
  \|v\|_V := \|\nabla v\|_{L^2(\Omega)}.
\]
(For $H^1_0(\Omega)$ this is equivalent to the usual $H^1$--norm by Poincaré.)
Define
\[
  a(u,v) := \int_\Omega \nabla u\cdot\nabla v\,dx,
  \qquad
  g(v) := \int_\Omega f\,v\,dx.
\]
Then the weak formulation of (strong) is indeed:
\[
  \boxed{
    \text{Find } u\in V \text{ such that }
    a(u,v) = g(v)\quad \forall v\in V.
  }
\]
\textbf{Boundedness of $g$.}
For $v\in V$,
\[
  |g(v)|
  = \left|\int_\Omega f\,v\,dx\right|
  \le \|f\|_{L^2(\Omega)}\,\|v\|_{L^2(\Omega)}.
\]
By the Poincaré inequality (given in the problem),
\[
  \|v\|_{L^2(\Omega)} \le C_\Omega \|\nabla v\|_{L^2(\Omega)} = C_\Omega \|v\|_V,
\]
so
\[
  |g(v)|
  \le C_\Omega \|f\|_{L^2(\Omega)}\,\|v\|_V,
\]
hence $g\in V^\ast$ is a bounded linear functional.
\subsection*{(c) Ellipticity in $H^1_0(\Omega)$ and symmetry of $a(\cdot,\cdot)$}
We show that $a$ is $V$--elliptic on $V=H^1_0(\Omega)$ and symmetric.
\textbf{Symmetry.}
For any $u,v\in V$,
\[
  a(u,v)
  = \int_\Omega \nabla u\cdot\nabla v\,dx
  = \int_\Omega \nabla v\cdot\nabla u\,dx
  = a(v,u),
\]
so $a$ is symmetric.
\textbf{Definition (V--elliptic).}
A bilinear form $a:V\times V\to\mathbb{R}$ is called $V$--elliptic (or coercive) if there
exists $\alpha>0$ such that
\[
  a(v,v) \ge \alpha\,\|v\|_V^2 \quad\forall v\in V.
\]
\textbf{Ellipticity on $H^1_0(\Omega)$.}
We have chosen the norm
\[
  \|v\|_V = \|\nabla v\|_{L^2(\Omega)}.
\]
Then, for $v\in V$,
\[
  a(v,v)
  = \int_\Omega |\nabla v|^2\,dx
  = \|\nabla v\|_{L^2(\Omega)}^2
  = \|v\|_V^2.
\]
Hence $a$ is $V$--elliptic with constant
\[
  \alpha = 1.
\]
(If we instead used the full $H^1$--norm, $a$ would still be elliptic by the Poincaré
inequality, with a different constant $\alpha>0$.)
Thus $a$ is symmetric and $H^1_0$--elliptic.
\subsection*{(d) Existence and uniqueness via Riesz (abstract Lax--Milgram)}
Let $V$ be a Hilbert space with inner product $\langle\cdot,\cdot\rangle_V$ and norm
$\|\cdot\|_V$.
Let $a:V\times V\to\mathbb{R}$ be a bilinear form satisfying:
\begin{itemize}
  \item \emph{Boundedness:}
        there exists $M>0$ such that
        \[
          |a(u,v)| \le M\,\|u\|_V\,\|v\|_V
          \quad\forall u,v\in V;
        \]
  \item \emph{Symmetry:}
        $a(u,v) = a(v,u)$ for all $u,v\in V$;
  \item \emph{V--ellipticity:}
        there exists $\alpha>0$ such that
        \[
          a(v,v) \ge \alpha\,\|v\|_V^2
          \quad\forall v\in V.
        \]
\end{itemize}
Let $g\in V^\ast$ be a bounded linear functional.
We want to prove that there is a unique $u\in V$ such that
\[
  a(u,v) = g(v)\quad\forall v\in V,
\]
i.e.\ that the abstract weak problem (2) is well posed.

\textbf{Step 1: Use Riesz representation to represent $g$.}
By the Riesz representation theorem (given in the problem), for the functional
$g\in V^\ast$ there exists a unique $w\in V$ such that
\[
  g(v) = \langle w,v\rangle_V \quad\forall v\in V.
\]
\textbf{Step 2: Associate $a(\cdot,\cdot)$ with a linear operator $A:V\to V$.}
For each fixed $u\in V$, the map
\[
  v \mapsto a(u,v)
\]
is a bounded linear functional on $V$ (by the boundedness assumption).
By Riesz, there exists a unique element $Au \in V$ such that
\[
  a(u,v) = \langle Au,v\rangle_V \quad\forall v\in V.
\]
This defines a linear operator $A:V\to V$ by assigning $u\mapsto Au$.


\emph{Properties of $A$:}
\begin{itemize}
  \item $A$ is linear and bounded, since
        \[
          \|Au\|_V
          = \sup_{\|v\|_V=1} |\langle Au,v\rangle_V|
          = \sup_{\|v\|_V=1} |a(u,v)|
          \le M \|u\|_V.
        \]
  \item $A$ is symmetric (self--adjoint w.r.t.\ $\langle\cdot,\cdot\rangle_V$), because
        \[
          \langle Au,v\rangle_V = a(u,v) = a(v,u) = \langle Av,u\rangle_V
          \quad\forall u,v\in V.
        \]
  \item Coercivity (ellipticity) translates into
        \[
          \langle Au,u\rangle_V = a(u,u) \ge \alpha\,\|u\|_V^2 \quad\forall u\in V.
        \]
\end{itemize}
\textbf{Step 3: Show $A$ is invertible.}
First, $A$ is injective: if $Au=0$, then
\[
  0 = \langle Au,u\rangle_V = a(u,u) \ge \alpha\|u\|_V^2.
\]
Thus $\|u\|_V=0$ and hence $u=0$.
Next, show that the range $\mathcal{R}(A)$ is all of $V$.
From coercivity we have, for all $u\in V$,
\[
  \alpha\|u\|_V^2 \le \langle Au,u\rangle_V \le \|Au\|_V\,\|u\|_V,
\]
so if $u\neq 0$,
\[
  \|Au\|_V \ge \alpha\|u\|_V.
\]
This implies that $A$ is bounded below:
\[
  \|Au\|_V \ge \alpha\|u\|_V \quad\forall u\in V.
\]
Hence $\mathcal{R}(A)$ is closed in $V$.
Consider the orthogonal complement of the range:
\[
  \mathcal{R}(A)^\perp
  = \{ z\in V : \langle Au,z\rangle_V = 0 \ \forall u\in V\}.
\]
For $z\in \mathcal{R}(A)^\perp$ we have
\[
  0 = \langle Au,z\rangle_V = a(u,z) \quad\forall u\in V.
\]
Since $a$ is symmetric, $a(u,z) = a(z,u)$, hence
\[
  a(z,u) = 0 \quad\forall u\in V.
\]
In particular, taking $u=z$ gives
\[
  a(z,z) = 0.
\]
By ellipticity,
\[
  0 = a(z,z) \ge \alpha\|z\|_V^2 \quad\Rightarrow\quad z=0.
\]
Thus $\mathcal{R}(A)^\perp = \{0\}$.
But in a Hilbert space we always have
\[
  \overline{\mathcal{R}(A)} = (\mathcal{R}(A)^\perp)^\perp.
\]
Since $\mathcal{R}(A)$ is closed and $\mathcal{R}(A)^\perp = \{0\}$, it follows that
\[
  \mathcal{R}(A) = V.
\]
Therefore $A:V\to V$ is bijective (injective and surjective), hence invertible and
$\|A^{-1}\| \le \alpha^{-1}$.

\textbf{Step 4: Solve the variational problem.}
We must find $u\in V$ such that
\[
  a(u,v) = g(v) \quad\forall v\in V.
\]
Using the Riesz representation $g(v) = \langle w,v\rangle_V$ and the operator identity
$a(u,v) = \langle Au,v\rangle_V$, this becomes
\[
  \langle Au,v\rangle_V = \langle w,v\rangle_V \quad\forall v\in V.
\]
By uniqueness of Riesz representatives, this is equivalent to
\[
  Au = w.
\]
Since $A$ is invertible, there exists a unique $u\in V$ solving $Au=w$, hence a unique
$u\in V$ solving
\[
  a(u,v) = g(v)\quad\forall v\in V.
\]
Thus: For any bounded, symmetric, V--elliptic bilinear form $a(\cdot,\cdot)$ and any $g \in V^\ast$,
\[
  \boxed{
    \text{ there is a unique } u\in V \text{ with } a(u,v)=g(v)\ \forall v \in V.
  }
\]

\textbf{Application to the weak form in (b).}
In part (b) we showed that
\[
  a(u,v) = \int_\Omega \nabla u\cdot\nabla v\,dx
\]
is symmetric and $V$--elliptic on $V=H^1_0(\Omega)$, and that
\[
  g(v) = \int_\Omega f v\,dx
\]
is a bounded linear functional on $V$. Therefore, by the abstract result just proved,
there exists a unique $u\in H^1_0(\Omega)$ solving
\[
  \int_\Omega \nabla u\cdot\nabla v\,dx
  = \int_\Omega f v\,dx
  \quad\forall v\in H^1_0(\Omega).
\]
Hence the weak formulation derived in (b) is \emph{well posed} (it has a unique solution
depending continuously on $f$).
